\begin{qfmethod}[Equity-index valuation for macro: top-down country and sector selection, price-to-book, dividend yield, P/E and forward P/E, free-cash-flow yield and market capitalisation to GDP, with Gliner's thresholds as rules of thumb]\label{meth:equity-index-valuation}
\qffield{Statement} Gliner describes a top-down approach to equities, from a global view to countries, then sectors, then optionally companies; if the investor knows where the next opportunity lies, the reward comes from concentrated rather than index exposure \citep[chap.~8, sec.\ ``Taking a Top-Down Macro Approach'', PDF pp.~215--216, Figure~8.1]{s7-gliner}.

His valuation measures (sec.\ ``Valuation Techniques for Equities'', PDF pp.~228--233) start with price-to-book, where ``a P/B greater than 2.5$\times$ is generally thought to be overbought, and a P/B lower than 1.5$\times$ is oversold'' (the S\&P 500 near 1.5 in 2008, p.~228). They continue with dividend yield, paid while waiting (p.~229); trailing and forward P/E, with earnings yield its reciprocal (pp.~230--232); price-to-free-cash-flow and its reciprocal, a Graham-and-Dodd screen (pp.~232--233); and stock market capitalisation to GDP, ``more worrisome'' the further above 100 percent and quickly rising, as when Japan's passed 150 percent in 1989 (p.~233).

\qffield{Notation} P/B price over book value per share; P/E price over earnings per share (forward: analysts' next-year estimate); FCF operating cash flow minus capital expenditure.

\qffield{Conditions} Gliner's thresholds as of 2013. The thresholds are rules of thumb for the S\&P 500's history, not levels that transfer across markets or eras: sector mix, accounting, intangible assets and interest rates move normal valuation ranges. Market capitalisation to GDP mixes a global firm's value with domestic output.

\qffield{Practice} For \nolinkurl{cross-asset-style-premia}, value in equity indices is measured cross-sectionally (cheap against expensive countries), not against a fixed threshold. Gliner's Tobin's $q$, which he says gives ``a score between 0 and 1'' (p.~228), is not bounded above by one: $q$ is market value over replacement cost and has exceeded one for long periods.

\qffield{Pitfalls} P/B above 2.5$\times$ as ``generally overbought'', as a timing rule, is untested; valuation predicts long-horizon returns weakly and timing poorly.

\qffield{See also} \hyperref[def:equity-indices-derivatives]{equity indices}; \hyperref[meth:equity-macro-leading-indicators]{equity leading indicators}; \hyperref[def:macro-regime-asset-map]{regime map}.


\end{qfmethod}
